\section{Rotation and a trace identity}\label{sec:rotation}

This section compares two ways of closing a positive translation path.
One closes it over $\mathbb F_q$; the other closes it after moving its
initial segment past Frobenius.  The comparison requires a connected
Picard group.  We construct that group using a divisible norm lift on the
splitting cover of the generic centralizer.  Its modulus retains the
entire level divisor, including its multiplicities.

\subsection{Levels, translations, and genericity}

Throughout this section $G$ may be split reductive.  Fix a split maximal
torus $T$, its Weyl group $W$, and the translation lattice
$\Lambda=X_*(T)$.  Put
\[
 \Lambda_0=\{a\in\Lambda:\langle\chi,a\rangle=0
                         \text{ for every }\chi\in X^*(G)\}.
\]
The positive chamber is the chamber of standard Bernstein translations
chosen in Section~\ref{sec:categorical}.  At distinct closed points
$x_1,\ldots,x_r$, outside the fixed level divisor $D$, impose Iwahori
level.  Write $d_j=\deg(x_j)$ and $q_j=q^{d_j}$.  Additional fixed
parahoric flags are allowed.  All modifications below preserve the full
trivialization along $D$.

If $G$ has a split connected center, quotient by a rational free lattice
$\Xi$ of central modifications whose degree map is injective and has
cocompact image in the free abelian degree lattice.  Such a lattice exists:
central cocharacters span the character-degree space over $\mathbb Q$,
and rational divisors on $X$ realize a finite-index subgroup of its degree
lattice.  Choose a representative of every resulting degree coset.
Every path used below has degree zero.  If it closes in the quotient,
its closing central modification has degree zero and is therefore the
identity.  The same holds for Frobenius closure, since the lattice is
rational and point Frobenius preserves character degrees.  Thus these
paths lift to the selected degree sectors before taking the quotient.
This observation also specifies the quotient version of all subsequent
constructions.

For a positive tuple $a=(a_j)$, set
\[
 L_a=\sum_j d_j\ell(t^{a_j}),\qquad
 T_a=\prod_j q_j^{-\ell(t^{a_j})/2}
                    1_{I_jt^{a_j}I_j}.
\]
Here the rational Iwahori has volume one.  The $T_a$ extend
multiplicatively to the Bernstein torus.  They are the
\emph{function-normalized} operators of Proposition~\ref{cat:hecke}.
In particular, at a closed point the character of the pure standard
kernel is $(-1)^{\ell(t^a)}T_a$, including when the geometric factors
are cyclically permuted by Frobenius.  We retain this correction for
arbitrary translation labels; we do not restrict the Bernstein algebra
to an even sublattice.

Choose $u_j$ in the closed positive chamber and in
$2|W|\Lambda_0$.  At the $d_j$ geometric points above $x_j$, repeat
$u_j$ identically.  Choose positive integral multiples $\varpi$ of the
semisimple fundamental weights that belong to $X^*(T)$ and are trivial
on the connected center.  Assume
\begin{equation}\label{rot:genericity}
 \sum_jd_j\langle\varpi,w_ju_j\rangle\ne0
       \qquad\text{for every such $\varpi$ and every }(w_j)\in W^r.
\end{equation}
When $G$ is a torus the condition is empty.  We do not require each
$u_j$ to be regular.

Write $\mathrm{Fr}$ for point Frobenius, reserving $F$ for the function
field.  Let $Y_u$ be the open translation shtuka: an object consists of a
bundle $\mathcal E$ with the specified levels and a meromorphic
isomorphism
\[
 \phi:\mathcal E\longrightarrow\mathrm{Fr}(\mathcal E)
\]
of exact Iwahori positions $t^{u_j}$ at the geometric legs, preserving
all levels away from them.  Equivalently one can keep the intermediate
bundles in the path.  Modifications at distinct legs commute, and the
open convolution maps used below identify these descriptions.  Put
\begin{equation}\label{rot:complex}
 C(u)=R\Gamma_c(Y_u,\mathbf k),\qquad
 M(u)=C(u)[L_u](L_u/2).
\end{equation}
All $\ell(t^{u_j})$ are even, so the shift and cyclic-permutation signs
in this normalization are trivial.  We use the natural Frobenius on
$M(u)$; in terms of $C(u)$ it is $q^{-L_u/2}\mathrm{Fr}^*$.
The complex $M(u)$ is the cyclic complex of the repeated standard
translation kernels $J_{u_j}$.  For a positive tuple $v$, their coherent
commutation with the kernels $J_{v_j}$ lets us apply the character
construction in Equation~\eqref{cat:character}.  Denote the resulting
function-normalized endomorphism of $M(u)$ by $T_v$, with the parity
correction specified above.

\begin{theorem}\label{rot:trace}
Let $u$ satisfy the chamber, divisibility, degree-zero, and genericity
conditions above.  Let $v$ be any fixed regular positive tuple in
$\Lambda_0^r$, without a parity restriction.  For all sufficiently
large integers $n$,
\begin{equation}\label{rot:normalized-trace}
 \sum_{[\mathcal E]}(T_{nu+v})_{\mathcal E,\mathcal E}
   =\operatorname{Tr}_{\mathrm{alt}}
           (\mathrm{Fr}^{n}T_v\mid M(u)).
\end{equation}
The diagonal sum uses ordinary matrix entries on the rational
isomorphism classes of bundles with level, and is finite.  The assertion
retains arbitrary multiplicities in $D$ and the allowed central quotient.
\end{theorem}

We first bound $Y_u$ and identify its translation operators with finite
correspondences.  The large-Frobenius trace formula then counts loops
whose rotation by the first $u$-segment is identified with Frobenius.
The ordinary
matrix trace instead counts rational loops.  The main comparison uses
a connected Picard group to identify these two descent groupoids;
the proof of Theorem~\ref{rot:trace} then assembles the counts.

\subsection{Boundedness and separation}

We first record the elementary translation facts needed to apply the
bundle bounds of Proposition~\ref{amp:geometry}.  If affine Weyl elements
$w,w'$ satisfy $\ell(ww')=\ell(w)+\ell(w')$, the open convolution map
\begin{equation}\label{rot:open-product}
 IwI\mathbin{\times^I}Iw'I\longrightarrow Iww'I
\end{equation}
is an isomorphism.  Indeed the root subgroups in the two inversion sets
are disjoint, and their ordered products give the root coordinates on
the inversion set of $ww'$.  Thus a path of exact position $ww'$ has a
unique intermediate flag of the prescribed type.  This statement, with
its intermediate isomorphisms, also holds for Hecke groupoids.  In
particular, translations in one closed chamber have unique
length-additive subdivisions.

For an ordinary local automorphism $g$ with
$g^m\in It^{mb}I$ for all $m\ge1$, its dominant Newton vector is $b$.
Here the Newton vector means the valuations of its semisimple part,
viewed as a cocharacter modulo $W$.  Here is a linear-algebra verification that also works on a wall.
In any highest-weight representation, an integral torus-weight lattice
is preserved by $I$.  The growth rates of the operator norms of $g^m$
and its dual are consequently the extremal pairings of its weights
with $mb$.  Over a finite splitting extension, put $g$ in Jordan form.
Its unipotent part has finite $p$-power order, and conjugation by the
fixed change-of-basis matrix changes logarithmic operator norms by a
bounded amount.  The same growth rates are therefore those of the
valuations of its semisimple eigenvalues.  As the highest weights vary,
these extremal pairings determine the dominant cocharacter, proving
that it is $b$.  No semisimplicity of the representation in positive
characteristic is needed.  If $g$ preserves a parabolic reduction, the valuation
of a character of that parabolic is its pairing with the Newton vector
in the corresponding Levi, hence with a Weyl conjugate of $b$ in $G$.

\begin{lemma}\label{rot:bounded}
Under Equation~\eqref{rot:genericity}, $Y_u$ is a separated
Deligne--Mumford stack of finite type, and $C(u)$ is a bounded complex
with finite-dimensional cohomology.  For a fixed regular positive
$v\in\Lambda_0^r$ and all sufficiently large integers $n$, the ordinary
rational loop groupoid of exact position $nu+v$ has finitely many
isomorphism classes and finite mass.  Only bundles of instability
$O(1+n)$ occur in it.
\end{lemma}

\begin{proof}
Suppose the instability in $Y_u$ were unbounded.  Proposition~\ref{amp:geometry}
would give a maximal-parabolic canonical reduction preserved by $\phi$,
where preservation means that the reduction is carried to its
Frobenius image.  Work first at a geometric point with all its data
defined over $\mathbb F_{q^b}$, enlarging $b$ to be divisible by every
$d_j$.  Iterating $\phi$ $b$ times gives an ordinary meromorphic
automorphism $\gamma$ of the bundle.  Equation~\eqref{rot:open-product}
puts all its powers, at every geometric point above $x_j$, in the
Iwahori cells of translations $mbu_j$.

Apply the preceding Newton calculation to the preserved reduction.
The valuation of its fundamental character at one geometric leg is
$b\langle\varpi,w_ju_j\rangle$ for some $w_j$.  The original
Frobenius path intertwines the consecutive local automorphisms and
their reductions.  Since the character of the split parabolic is also
preserved, its valuation is the same at all $d_j$ geometric points
above $x_j$.  The product formula for the meromorphic character
therefore gives
\[
 b\sum_jd_j\langle\varpi,w_ju_j\rangle=0,
\]
contrary to Equation~\eqref{rot:genericity}.  Proposition~\ref{amp:geometry}
now bounds the instability, so boundedness of bundles and of the Hecke
positions gives finite type.  This argument includes wall labels, which
merely allow more Weyl choices.

The shtuka condition makes inertia unramified.  Infinitesimally an
automorphism compatible with $\phi$ equals its transported Frobenius
pullback; the latter has zero differential, so its tangent vector is
zero.  Automorphism schemes of a bundle with bounded level are affine
of finite type.  Thus the inertia is quasi-finite, and $Y_u$ is
Deligne--Mumford.

For separation, let two objects over a discrete valuation ring $R$ be
isomorphic over its fraction field.  Choose a faithful representation
of $G$.  At the generic point of the special fiber of $X_R$, both path
maps are integral isomorphisms: the legs are horizontal.  If $m$ is the
smallest valuation of a matrix entry of the generic isomorphism,
compatibility with the paths gives $m=qm$, since base Frobenius
multiplies vertical valuations by $q$ and multiplication by an integral
invertible matrix preserves the smallest valuation.  Hence $m=0$.
The same argument applies to the inverse.  The isomorphism consequently
extends at that vertical generic point.  All horizontal codimension-one
points are already covered by the generic-fiber bundle isomorphism.
Normality of the smooth surface $X_R$ and affineness of the isomorphism
torsor extend it across codimension two.  The path identities extend by
density.  So do the prescribed level identities: $D_R$ is flat over
$R$, including when $D$ is nonreduced.  The same argument treats fixed
parahoric flags.  This proves the valuative criterion for the
quasi-finite diagonal, which is therefore proper.

For completeness, finite-dimensional compact cohomology can be reduced
to the corresponding assertion for algebraic spaces.  On this bounded
family, choose a sufficiently large effective divisor $N$ disjoint
from the legs and the existing level such that an automorphism equal to the identity on $N$
is the identity.  A faithful representation reduces this to the
uniform vanishing of sections of a bounded family of endomorphism
bundles twisted by $-N$.  Frobenius-compatible trivializations on $N$
form a finite \'etale torsor under
\[
 H_N=\bigl(\operatorname{Res}_{N/\mathbb F_q}(G_N)\bigr)(\mathbb F_q).
\]
Indeed the jet group $\operatorname{Res}_{N/\mathbb F_q}(G_N)$ is smooth
and connected, and its Lang map is finite \'etale and surjective.
The resulting rigidified shtuka is a separated algebraic space of
finite type with a finite-group quotient equal to $Y_u$.  Compact
cohomology is therefore bounded and finite-dimensional; taking
$H_N$-invariants is exact over $\mathbf k$.  The same reduction applies
to weights.  A finite-type algebraic space over a field admits a finite
stratification by schemes, obtained by successively choosing a dense
open subspace that is a scheme.  Excision transfers the scheme weight
bounds to this algebraic space, and exact finite-group invariants retain
them.  Thus the pure constant coefficient in
Equation~\eqref{rot:complex} has the usual integral-weight compact
cohomology bounds, without any appeal to a weight theorem for arbitrary
Artin stacks.

Finally, $a=nu+v$ is regular positive, and for each fixed pairing in
Equation~\eqref{rot:genericity} the corresponding pairing for $a$ is a
nonconstant affine function of $n$.  Thus $a$ satisfies the same
nonvanishing condition for all sufficiently large $n$.  A rational
ordinary loop of type $a$ with excessive instability has, by
Proposition~\ref{amp:geometry}, a preserved maximal-parabolic reduction.
The same character and product-formula argument excludes it.  The
quantitative bound is $O(1+n)$.  There are finitely many rational bundles
in that range, and each admits only finitely many rational
meromorphic automorphisms with these pole bounds: in a faithful
representation they lie in a fixed finite-dimensional space of
sections over a finite field.  The automorphism group of every such
rational loop has finitely many rational points.  This proves the
last assertion.
\end{proof}

\subsection{The finite correspondence underlying a translation}

For a positive tuple $v$, let $\mathcal C_v$ classify a commuting square
\[
\begin{tikzcd}[column sep=large,row sep=large]
 \mathcal E \arrow[r,"h", "v"'] \arrow[d,"\phi"'] &
 \mathcal E' \arrow[d,"\phi'"] \\
 \mathrm{Fr}(\mathcal E) \arrow[r,"\mathrm{Fr}(h)"'] &
 \mathrm{Fr}(\mathcal E').
\end{tikzcd}
\]
The vertical paths have type $u$, and $h$ has type $v$.  Denote the
initial and final shtuka projections by $p$ and $q$ respectively.
The direction used on cohomology is $p_!q^*$, pullback towards the
initial vertex.

\begin{lemma}\label{rot:finite-correspondence}
The two projections of $\mathcal C_v$ are finite \'etale after passing
to perfections.  On $C(u)$ the function-normalized categorical
translation is
\begin{equation}\label{rot:geometric-operator}
 T_v=q^{-L_v/2}p_!q^*.
\end{equation}
These correspondences compose according to addition in a closed
chamber.  For $v=u$ the correspondence is the graph of point
Frobenius, so $T_u$ acts as Frobenius on $M(u)$.
\end{lemma}

\begin{proof}
We first work at one geometric leg.  Take a reduced gallery for $t^v$.
After a prefix $b$, cyclic rotation changes the shtuka label to the
translation $\tau=b^{-1}t^ub$.  Since $u$ and $v$ lie in one closed
chamber,
\[
 \ell(t^ub)=\ell(t^u)+\ell(b).
\]
For the next simple reflection $s$, one has
$\ell(\tau s)=\ell(\tau)+1$.  A translation has opposite left and right
length changes at a wall to which its vector is not parallel.
Consequently either $s$ is an initial letter of $\tau$, or $\tau$
is parallel to the $s$-wall and commutes with $s$.

In the first case write $\tau=sw$.  Unique open subdivision identifies
the step with the cyclic rotation from the path $sw$ to the path $ws$.
It is an isomorphism on perfections.  The cancellation of the inverse
simple kernel against that initial factor in the categorical action
is precisely the cyclic isomorphism of Proposition~\ref{cat:cyclic}.
Thus its unnormalized action is pullback along this geometric step.

For a closed leg of degree $d$, perform this rotation on the whole block
of $d$ geometric factors.  Cyclicity applied to
$\Delta_s^{\boxtimes d}\star\Delta_w^{\boxtimes d}$ changes the repeated
label $\tau=sw$ to the repeated label $ws=s\tau s$; unique reduced
subdivisions identify both ends with the stated path stacks.

In the parallel case, pass from the alcove to its $s$-panel.  The partial
relative position identifies the two rank-one flag varieties: on the
root groups for that panel, translation conjugation has exponent zero.
A refinement of the partial shtuka is therefore a flag fixed by the
resulting Frobenius-semilinear identification.  Over a degree-$d$ closed
leg, these refinements form a finite \'etale family of $q^d+1$ flags.
To check this over a general base, trivialize the rank-one torsor
\'etale locally and use the Lang map for its connected rank-one group;
the fixed flags then become $\mathbb P^1(\mathbb F_{q^d})$.
If $\pi$ denotes the map forgetting this refinement, the closed simple
correspondence acts by $\pi^*\pi_!$, and the open simple correspondence
acts by
\begin{equation}\label{rot:panel-action}
 \pi^*\pi_!-\operatorname{id}.
\end{equation}

Here Equation~\eqref{rot:panel-action} is also an equality with the
categorical transfer, rather than only an equality of trace functions.
Use the unshifted closed rank-one kernel, which factors through the
panel projection.  Its unit and counit give the diagonal and projection
maps for the proper flag variety.  After restriction to the twisted
trace, the Frobenius graph meets the diagonal transversely: its
linearized fixed-point equation has derivative $-1$.  Every fixed flag
therefore has coefficient one in the unit/counit transfer.  The
projection formula identifies the resulting operator with sum and
pullback over the finite family of fixed flags.  This is the explicit
rank-one calculation in Proposition~\ref{cat:hecke}.  Subtracting the
identity kernel by the standard open/closed triangle gives
Equation~\eqref{rot:panel-action}.  In particular no compact cohomology
shift of the flag variety remains in this finite transfer.

At a closed leg the same calculation is made in the product of its
$d$ geometric flag varieties, with Frobenius cyclically permuting the
factors.  A fixed tuple is determined by one
$\mathbb F_{q^d}$-flag.  The same transverse intersection proves the
coefficient-one assertion and gives the parameter $q^d$.
Graded cyclic permutation produces the single-factor parity in the
pure standard kernel.  The correction defining $T_v$ removes that
parity even when $v$ is odd.  Thus normalization gives the factor
$q^{-L_v/2}$ in Equation~\eqref{rot:geometric-operator}.

Length-zero gallery steps identify the relevant paths.  Composing the
two kinds of simple steps proves finiteness on perfections and the
operator assertion.  Unique subdivision also proves compatibility of
composition and of commuting-translation shuffles.  When $v=u$, the
two initial steps in the commuting square are subdivisions of the same
path of position $2u$ and hence coincide.  The square is therefore
$h=\phi$, $\mathcal E'=\mathrm{Fr}(\mathcal E)$, and
$\phi'=\mathrm{Fr}(\phi)$.  Its action is $\mathrm{Fr}^*$ on $C(u)$;
after the normalization in Equation~\eqref{rot:complex} it is
Frobenius on $M(u)$.
\end{proof}

\subsection{The Frobenius trace and the shifted loop condition}

We use the following precise form of the large-Frobenius trace theorem.
If $A\xleftarrow{c_2}C\xrightarrow{c_1}A$ is a correspondence of
separated finite-type schemes over $\mathbb F_q$, with $c_1$ proper and
$c_2$ quasi-finite, then, for all sufficiently large $n$, its
Frobenius-twisted fixed-point set is finite and the alternating compact
trace equals the sum of the naive local traces.
This is \cite[Theorem~2.3.2(a),(b) and Remark~2.3.3(c)]{Varshavsky},
with the open subset there equal to all of $A$.  The same proof applies
to compactifiable algebraic spaces, as stated in the final paragraph
of that paper's introduction.  In our application both projections
are finite and the coefficient is constant, so each point counted by
the pullback/sum correspondence has local trace one.

Here is the reduction that avoids invoking such a theorem directly
for a stack.  Use the rigidifying divisor $N$ from the proof of
Lemma~\ref{rot:bounded}.  Every Hecke step at the legs transports its
$N$-trivialization uniquely.  Hence the correspondence lifts to the
rigidified algebraic space $\widetilde Y_u$, commutes with the constant
finite group $H_N$, and has finite projections.  Exact averaging gives
\begin{equation}\label{rot:averaging}
 \operatorname{Tr}_{\mathrm{alt}}(\mathrm{Fr}^{n*}U_v\mid C(u))
 =\frac1{|H_N|}\sum_{h\in H_N}
   \operatorname{Tr}_{\mathrm{alt}}
       (h\,\mathrm{Fr}^{n*}\widetilde U_v\mid
                     R\Gamma_c(\widetilde Y_u,\mathbf k)),
\end{equation}
where $U_v=p_!q^*$ is unnormalized.  Apply the theorem to the finitely
many lifted correspondences.  The fixed-point groupoid downstairs is
the quotient of the disjoint union of their twisted fixed points by
$H_N$, with $H_N$ acting by conjugation on the twisting element.
Its weighted cardinality is exactly the right side of
Equation~\eqref{rot:averaging} after replacing traces by counts.
This proves the stack formula with all automorphism factors.

If a gallery rotation uses an inverse purely inseparable map, choose a
fixed Frobenius power that clears its denominator, and do the same for
the finitely many maps in the correspondence.  The resulting
correspondence has finite algebraic models.  Composing it with
$\mathrm{Fr}^{n-b}$ gives the original correspondence composed with
$\mathrm{Fr}^n$.  Since compact \'etale cohomology and fixed geometric
points are invariant under the intervening universal homeomorphisms,
the preceding argument applies once $n$ exceeds this fixed $b$.
The Frobenius operators in this procedure are transported with the
correspondence; none is replaced by the identity.

Let $\mathcal L_a$ denote the geometric groupoid of ordinary meromorphic
loops of exact position $a$: its objects are $(\mathcal E,\gamma)$,
where $\gamma$ is a self-modification preserving the levels.  For
$a=nu+v$, let $R_u$ rotate the first $u$-segment to the end.  It is
invertible when $a$ is regular, as will also follow from the torus
apartment description below.  Write
\[
 \mathcal L_a^{\mathrm{sh}}
   =\{(\mathcal E,\gamma,\iota):
           \iota:R_u(\mathcal E,\gamma)
                    \xrightarrow{\sim}\mathrm{Fr}(\mathcal E,\gamma)\}.
\]
Arrows in this notation must respect $\iota$.

\begin{lemma}\label{rot:concatenation}
For sufficiently large $n$, the fixed-point groupoid for
$\mathrm{Fr}^{n*}U_v$ on $Y_u$ is equivalent to
$\mathcal L_{nu+v}^{\mathrm{sh}}$.  Consequently
\begin{equation}\label{rot:shifted-trace}
 \operatorname{Tr}_{\mathrm{alt}}(\mathrm{Fr}^{n*}U_v\mid C(u))
                  =|\mathcal L_{nu+v}^{\mathrm{sh}}|.
\end{equation}
\end{lemma}

\begin{proof}
With the direction $p_!q^*$ used above, the fixed-point identification
is $p=\mathrm{Fr}^nq$.  Thus its data consist of a shtuka
$(\mathcal E,\phi)$ and a closing step
$h:\mathrm{Fr}^n\mathcal E\to\mathcal E$ of type $v$, compatible with
$\phi$.  Concatenate
\[
 \mathcal E\xrightarrow{\phi}\mathrm{Fr}\mathcal E
       \longrightarrow\cdots\longrightarrow
 \mathrm{Fr}^n\mathcal E\xrightarrow{h}\mathcal E.
\]
The total position is $nu+v$ by Equation~\eqref{rot:open-product}.
Compatibility with $h$ identifies rotation of the initial $u$-step
with point Frobenius.

Conversely, subdivide a loop of position $nu+v$ into $n$ successive
$u$-steps and one $v$-step.  The given identification of the first
rotation with Frobenius turns the first step into
$\phi:\mathcal E\to\mathrm{Fr}\mathcal E$; subsequent steps are its
Frobenius translates.  If $\gamma$ is the total loop, the last step is
$h=\gamma(\phi^{(n)})^{-1}$.  The identity
$\phi\gamma=\mathrm{Fr}(\gamma)\phi$ gives the required commuting
square for $h$.  Unique open subdivision makes these constructions
inverse on objects, intermediate isomorphisms, and arrows.  The trace
formula just proved gives Equation~\eqref{rot:shifted-trace}.
\end{proof}

\subsection{Regular centralizers and connected norm lifts}

We will compare the shifted count with the ordinary rational loop
count.  Only their fixed-point groupoids need be finite; the geometric
groupoid $\mathcal L_a$ itself need not be bounded.

For a torus $G$, the centralizer is already split and all the
degree-zero labels are zero, so the following apartment assertion is
empty.  Otherwise, for a regular positive $a_j$, an ordinary local automorphism $g$ of
position $t^{a_j}$ has all powers in the positions $t^{ma_j}$ by
Equation~\eqref{rot:open-product}.  We explain why this forces both
strong regularity and the apartment assertion needed for rotation.
Work over a finite extension over which the object and its local flag
are defined, and use the reduced affine building of $G$ there.  Let
$A_0$ be the base alcove.  Its Iwahori fixes it pointwise
\cite[Sections~1.2 and~3.7]{TitsReductiveLocal}. For every
point $z\in A_0$, writing $g^m=i_1t^{ma_j}i_2$ gives
\[
 d(z,g^mz)=d(z,t^{ma_j}z)=m\|a_j\|.
\]
Equality in the triangle inequality for $m=2$ shows that the segments
$[g^kz,g^{k+1}z]$ concatenate to a geodesic axis.  Its direction is
regular, since its vector distance is $a_j$.  After taking a power to
remove the finite permutation of building types, the minimum set of
this regular-axis isometry is a unique apartment; this is
\cite[Lemma~2.3]{CapraceCiobotaru}.  The original $g$ preserves this
apartment because it commutes with its power.  Its finite Weyl part
fixes the regular translation vector and hence is trivial.  Thus $g$
acts on the apartment by a translation. An apartment corresponds to a
maximal split torus $S$, and its stabilizer is $N_G(S)$
\cite[Section~2.1]{TitsReductiveLocal}. The trivial finite Weyl part
therefore puts $g$ in $Z_G(S)=S$, since $G$ is split. The central
Euclidean factor for reductive $G$ can simply be restored; the labels
under consideration have zero character degree.

Every point of $A_0$ lies on such an axis, so the original flagged
lattice belongs to that torus apartment.  The valuations of its root
characters are the root pairings with $a_j$, all nonzero.  In particular
no root value is one.  Moreover, a Weyl element centralizing $g$ would
fix its regular valuation vector and must be the identity.  Hence $g$
is strongly regular semisimple and this split torus is its full
centralizer.  Unique subdivision in the apartment now
identifies $R_u$ with local modification by the ordered cocharacter
$u_j$.  Modification by $-u_j$ is its inverse.  This description is
unchanged by further extension of the residue field.

The generic conjugacy invariant of a loop fixed by either ordinary or
shifted Frobenius is rational: rotation does not change that invariant.
Fix one such invariant $c$.  It determines a canonical torus $S_c/F$
with a tautological element $\gamma_c\in S_c(F)$.  This construction is
obtained by descending the split torus along the Weyl torsor above $c$.
For any realization of $c$ as a strongly regular automorphism, its
centralizer is canonically identified with $S_c$: two conjugating
choices differ by an element of the same commutative centralizer.
The preceding argument first splits this torus after a finite residue
extension.  Over the original completed field at a closed leg, local
Galois acts on its split labels through $W$ and fixes the valuation
vector of the rational tautological element.  That vector is regular,
so its Weyl stabilizer is trivial.  Thus $S_c$ is already split over the
original completed field.  The valuations distinguish the ordering,
so the rotation displacements are rational there as well.

\begin{lemma}\label{rot:anisotropic}
If $a$ satisfies Equation~\eqref{rot:genericity}, with $a$ replacing $u$,
then $S_c$ has no split cocharacter directions outside the connected
center of $G$.
\end{lemma}

\begin{proof}
A noncentral invariant cocharacter can be placed in the closure of a
Weyl chamber.  Its stabilizer is the Weyl group of a proper standard
Levi.  Choose a fundamental character for a simple root missing from
that Levi, multiplied to be integral.  It is fixed by the arithmetic
monodromy of $S_c$, and so defines a rational character $\chi$ of that
torus.  The value $\chi(\gamma_c)$ is a rational function on $X$.
Away from the legs it is a unit: any realization of the loop is an
integral automorphism there, as is its inverse, and hence all its
torus-character valuations vanish.  At a leg its valuation is
$\langle\varpi,w_ja_j\rangle$ for a Weyl choice $w_j$.  The product
formula contradicts the assumed nonvanishing.  This argument uses
the tautological torus element and therefore applies even before a
rational bundle realizing $c$ has been chosen.
\end{proof}

\begin{lemma}\label{rot:picard}
Suppose $a=nu+v$ is regular and satisfies the nonvanishing condition,
and suppose the two fixed-point groupoids in question are finite.
Then ordinary and shifted Frobenius descent on $\mathcal L_a$ have
the same weighted count:
\begin{equation}\label{rot:descent-count}
 |\mathcal L_a(\mathbb F_q)|=|\mathcal L_a^{\mathrm{sh}}|.
\end{equation}
The assertion holds with arbitrary multiplicities in $D$ and after
the central quotient described above.
\end{lemma}

\begin{proof}
We give the comparison separately for each rational invariant $c$.
First we construct a connected group acting on suitable full-level
subgroupoids of loops.  We then realize rotation as a rational point
of that group and use Lang's theorem to compare the two descent
conditions.
Let $Y\to X$ be the normalization in a connected arithmetic Galois
splitting field of $S_c$.  Its group $\Gamma$ is a subgroup of $W$;
write $N=|\Gamma|$.  This cover need not be geometrically connected.
It is split and unramified at the regular legs.  Let
$\mathsf M=X_*(T)$ be the abstract split cocharacter lattice used on
the cover.  There is a norm homomorphism of generic tori
\begin{equation}\label{rot:norm}
 \nu:\operatorname{Res}_{F(Y)/F}
          (\mathbb G_m\otimes_{\mathbb Z}\mathsf M)\longrightarrow S_c,
\end{equation}
which multiplies Galois conjugates with their descent action on
$\mathsf M$.  We will use a connected source for this norm.  Neither
surjectivity nor separability of the norm itself is required.

Choose a finite rational set $Z\subset X$ disjoint from the legs and
containing the support of $D$, the ramification points, the points
where $\gamma_c$ is not regular integral, and all other prescribed
level points.  Add one auxiliary closed point outside the legs.
For an effective rational modulus $B$ supported on the full inverse
image of $Z$, with positive multiplicity at every such point, put
\begin{equation}\label{rot:picard-source}
 Q_B=\operatorname{Pic}^{\underline 0}(Y,B)
                       \otimes_{\mathbb Z}\mathsf M.
\end{equation}
The superscript means degree zero on \emph{each geometric component}.
The modulus meets every component, so a line bundle with its specified
$B$-trivialization has no scalar automorphisms.  Consequently $Q_B$ is
an actual smooth connected commutative algebraic group.  One can see
this from the generalized-Picard exact sequence: over the algebraic
closure it is an extension of the product of the component Jacobians
by the products of the local truncated unit groups, modulo one
constant multiplicative group per component.  A truncated unit group
is smooth and connected, including for nonreduced moduli in positive
characteristic.  This is also the precise representability and
connectedness statement in \cite[Section~1.4, Proposition~1.6 and the
following paragraph]{JordanRibetScholl}.  The auxiliary point rigidifies this source
only; it adds no level to the original loop problem.

For the descent comparison, it suffices to construct a coherent action
of $Q_B(\overline{\mathbb F}_q)$ on geometric objects and arrows,
compatible with Frobenius.  We now impose the integral conditions that
retain the full level.  Let $\mathcal L_{a,c,B}$ be
the full subgroupoid on objects for which norms of split congruence
units at $B$ preserve the local lattices, all flags, and the complete
$D$-trivialization.  Every object belongs to such a subgroupoid after
increasing the multiplicities of $B$.  Indeed the generic torus
embedding and the norm are continuous for the local valuation
topologies; sufficiently deep units map into the prescribed principal
congruence stabilizer.  Away from $Z$ and the legs the regular integral
centralizer is an unramified torus, and its integral units preserve the
lattice.  At the legs this follows from the centralizer-apartment
description above.

Represent an element of $Q_B(\overline{\mathbb F}_q)$ by an
$\mathsf M$-valued divisor disjoint from $B$.  Such a representative
exists by choosing meromorphic
sections agreeing with its trivializations to the specified finite
orders.  Modify the bundle at the image of its support by the norm
cocharacters in Equation~\eqref{rot:norm}: change the local
trivializations at those points and glue them to the unchanged torsor
away from the support.  This uses formal gluing for smooth affine
torsors on a curve, applied at finitely many points; see
\cite[arXiv version, Section~5.a, Lemma~5.1]{PappasZhu}.
At $Z$ keep the original lattice and its level.
All modifications commute with $\gamma_c$, so the result remains a
loop with the same invariant and the same relative positions at the
legs.  If two representatives differ by the divisor of a rational
split-torus function $h$ with $h|_B=1$, the meromorphic automorphism
$\nu(h)$ identifies the two results.  At $Z$ it extends and preserves
the full nonreduced level by the defining stabilizer condition;
elsewhere it is the usual change of lattice trivialization.
The normalized function $h$ is unique: the ratio of two choices is
constant on each geometric component and equals one along $B$.
The ratios of three representatives multiply, so these isomorphisms
satisfy their cocycle identities.  Tensor product therefore gives a
coherent action of $Q_B(\overline{\mathbb F}_q)$ on the groupoid.
An arrow between loops identifies their canonical centralizers and
carries each norm modification to the corresponding one.  Thus the
action is natural in all arrows, including every target automorphism.

The action preserves $\mathcal L_{a,c,B}$.  Namely a torus modification
does not change which elements of the commuting norm-unit subgroup
stabilize a lattice.  The same observation proves stability under
$R_u$.  All these constructions are Frobenius-compatible, since the
cover and the modulus are rational.

We next exhibit rotation in the connected source
Equation~\eqref{rot:picard-source}. At a geometric point $y$ of $Y$
over the leg $x_j$, let $\nu_y\in\mathsf M\otimes\mathbb Q$ be the
valuation vector of the tautological element $\gamma_c$ in the split
coordinates on $Y$. It belongs to the Weyl orbit of $a_j$.
Since $a_j$ is regular, there is a unique $w_y\in W$ with
$\nu_y=w_ya_j$.  Define the $\mathsf M$-valued divisor
\[
 \Delta=\sum_j\ \sum_{y\mapsto x_j}w_y(u_j/N)[y],
\]
where the inner sum runs over all geometric points above $x_j$.
These cocharacters are integral because $u_j\in2|W|\Lambda_0$ and
$N$ divides $|W|$. The valuation vectors are transported by the deck
action and preserved by Frobenius on the abstract constant lattice
$\mathsf M$. Uniqueness of $w_y$ therefore makes $\Delta$ both rational
and $\Gamma$-equivariant with its descent action on the labels.
Under the ordered local identification, every lift contributes
$u_j/N$ to its norm.  There are $N$ lifts over each geometric leg, so
the norm is exactly the rotation modification by $u_j$.
The support of $\Delta$ avoids $B$, giving its canonical
trivialization along the full modulus.

We verify its multidegree.  Let $\Gamma_{\mathrm{geom}}$ be the
geometric monodromy subgroup, and write $e_C$ for the vector degree of
$\Delta$ on a geometric component $C$.  Deck equivariance gives
\[
 e_{\gamma C}=\gamma e_C,
\]
so $\Gamma_{\mathrm{geom}}$ fixes $e_C$.  Frobenius preserves the labels
of the abstract split torus, and rationality gives
$e_{\mathrm{Fr}C}=e_C$.  Choose $\sigma\in\Gamma$ inducing the same
permutation of geometric components as Frobenius.  Then
\[
 \sigma e_C=e_{\sigma C}=e_{\mathrm{Fr}C}=e_C.
\]
The coset of $\sigma$ generates the cyclic group
$\Gamma/\Gamma_{\mathrm{geom}}$.  Together with invariance under
$\Gamma_{\mathrm{geom}}$, this proves
\begin{equation}\label{rot:component-degree}
 e_C\in(\mathsf M\otimes\mathbb Q)^\Gamma.
\end{equation}
By Lemma~\ref{rot:anisotropic}, this vector is central.  Every coefficient
of $\Delta$ lies in the rational subspace annihilating $X^*(G)$, and so
does $e_C$.  That derived subspace meets the rational central
cocharacter subspace in zero.  Thus $e_C=0$ on every component.
This is the point at which arithmetic, rather than merely geometric,
monodromy is needed.  We have obtained an element
$a_B\in Q_B(\mathbb F_q)$ whose action is $R_u$.

We can now compare descent.  Lang's theorem applies to the smooth
connected source $Q_B$; see \cite[Theorem~1 and its corollary]{Lang}.  In this
commutative case its surjectivity is especially direct:
$\mathrm{Fr}-1$ has derivative $-1$, hence has open image, and an open
subgroup of a connected algebraic group is the entire group.  Choose
$b\in Q_B(\overline{\mathbb F}_q)$ such that
\[
 \mathrm{Fr}(b)-b=a_B.
\]
If $\mathsf t_b$ denotes its action, Frobenius compatibility supplies
a natural identification
\[
 R_u^{-1}\mathrm{Fr}\,\mathsf t_b
                    \simeq\mathsf t_b\,\mathrm{Fr}.
\]
Hence $\mathsf t_b$, with inverse $\mathsf t_{-b}$, gives an equivalence
of descent groupoids
\begin{equation}\label{rot:lang-equivalence}
 \operatorname{Fix}(\mathrm{Fr},\mathcal L_{a,c,B})
   \simeq
 \operatorname{Fix}(R_u^{-1}\mathrm{Fr},\mathcal L_{a,c,B}).
\end{equation}
The identity in the source is an isomorphism of rigidified line
bundles, with no scalar ambiguity.  Equation~\eqref{rot:lang-equivalence}
therefore includes the Frobenius isomorphisms and their compatible
arrows.  It preserves automorphism groups, not only geometric
isomorphism classes.

We also spell out effectivity of the descent used here.  An object and
a Frobenius isomorphism are defined over a finite extension.  After
iterating by that extension degree, the resulting automorphism belongs
to the rational points of an affine finite-type group over a finite
field, hence has finite order.  A further iterate is the identity.
The datum is thus finite Galois descent, which is effective for bundles,
levels, and their meromorphic isomorphisms.  Ordinary Frobenius descent
is precisely the rational-point groupoid in
Equation~\eqref{rot:descent-count}.  The shifted version can either be
read from Equation~\eqref{rot:lang-equivalence} or from the explicit
concatenation construction of Lemma~\ref{rot:concatenation}.

The subgroupoids $\mathcal L_{a,c,B}$ exhaust the geometric loop
groupoid as the multiplicities grow.  Both fixed-point groupoids have
finitely many isomorphism classes by hypothesis.  Choose $B$ large
enough to contain representatives of both.  The equivalence then
proves equality of their entire weighted counts for $c$.  Summing over
the finitely many contributing rational invariants proves
Equation~\eqref{rot:descent-count}.

Finally, a norm modification arising from componentwise degree-zero
source line bundles changes no character degree of $G$.  Its action
commutes with central modifications and therefore descends to the
quotient by $\Xi$.  The initial degree-lifting observation identifies
its closed paths with the ones just considered.  This proves the
central-quotient assertion as well.
\end{proof}

\subsection{Proof of the trace identity}

\begin{proof}[Proof of Theorem~\ref{rot:trace}]
With $U_v=p_!q^*$ on the unnormalized complex, the identity to prove is
equivalently
\begin{equation}\label{rot:raw-trace}
 \sum_{[\mathcal E]}(T_{nu+v})_{\mathcal E,\mathcal E}
  =q^{-(nL_u+L_v)/2}
       \operatorname{Tr}_{\mathrm{alt}}
           (\mathrm{Fr}^{n*}U_v\mid C(u)).
\end{equation}
Set $a=nu+v$.  For $n$ large it is regular and satisfies the
nonvanishing condition.  Lemma~\ref{rot:bounded} makes the ordinary
loop count finite, and the large-Frobenius formula makes the shifted
count finite.  For a fixed rational bundle $\mathcal E$, closing a
Hecke modification whose target is isomorphic to $\mathcal E$ amounts
to choosing such an isomorphism, modulo its automorphisms.  Consequently
its diagonal matrix entry is the corresponding loop mass.  Summing
gives
\[
 \begin{aligned}
 \sum_{[\mathcal E]}(T_a)_{\mathcal E,\mathcal E}
       &=q^{-L_a/2}|\mathcal L_a(\mathbb F_q)|\\
       &=q^{-L_a/2}|\mathcal L_a^{\mathrm{sh}}|\\
       &=q^{-L_a/2}\operatorname{Tr}_{\mathrm{alt}}
          (\mathrm{Fr}^{n*}U_v\mid C(u)).
 \end{aligned}
\]
The first equality includes the stack multiplicities; no extra reciprocal
automorphism factor is to be inserted in the matrix trace.
The second equality is Lemma~\ref{rot:picard}, and the third is
Lemma~\ref{rot:concatenation}.  Since lengths add in the positive chamber,
$L_a=nL_u+L_v$, proving Equation~\eqref{rot:raw-trace}.
Equation~\eqref{rot:geometric-operator} and the even shift in
Equation~\eqref{rot:complex} give
Equation~\eqref{rot:normalized-trace}.  The operator $T_v$ here is
function-normalized; the parity correction for an arbitrary $v$ was
kept in Lemma~\ref{rot:finite-correspondence}.
\end{proof}
